\section{Метод решения}

Логика программы разделена на три независимых уровня:
\begin{enumerate}
  \item \emph{общий интерфейс} \texttt{src/os\_api.h} --- абстракции \texttt{os\_pipe} и
    \texttt{os\_process}, а также единый набор функций для создания и использования каналов и
    процессов;
  \item \emph{платформозависимые реализации} \texttt{src/os\_api\_win.c} (Win32 API) и
    \texttt{src/os\_api\_posix.c} (POSIX API) --- здесь находятся все системные вызовы, условная
    компиляция отсутствует, выбор файла выполняется в \texttt{CMakeLists.txt};
  \item \emph{прикладной уровень} \texttt{src/app/parent.c} и \texttt{src/app/child.c} --- вызывают только
    функции единого API, никаких системных вызовов в коде приложений нет.
\end{enumerate}

Алгоритм работы родительского процесса:
\begin{enumerate}
  \item создаются два канала: \texttt{in} (родитель $\rightarrow$ ребёнок) и \texttt{out}
    (ребёнок $\rightarrow$ родитель);
  \item через \texttt{os\_process\_spawn} запускается \texttt{child} с именем файла-отчёта;
    обёртка самостоятельно подключает каналы к стандартному вводу/выводу дочернего процесса и
    закрывает ненужные в текущем процессе концы каналов;
  \item в цикле читаются строки со стандартного ввода, числа с плавающей запятой извлекаются и
    записываются в канал \texttt{in};
  \item по окончании ввода закрывается записывающий конец канала \texttt{in} --- это сигнал
    \texttt{EOF} для дочернего процесса;
  \item родитель читает итоговую сумму из канала \texttt{out}, закрывает каналы, ожидает
    завершения \texttt{child} (\texttt{os\_process\_wait}) и закрывает дескриптор процесса.
\end{enumerate}

Дочерний процесс получает числа через стандартный ввод, суммирует их, записывает строку отчёта в
заданный файл и передаёт сумму обратно через стандартный вывод. Выбор низкоуровневых деталей
(преобразование argv в командную строку Windows, добавление \texttt{.exe}, префикс \texttt{./}
на POSIX и т.\,п.) полностью инкапсулирован в реализации обёртки.

\section{Описание программы}

\textbf{Разделение по файлам:}

\begin{center}
\begin{tabular}{ll}
\texttt{src/os\_api.h}       & единый интерфейс (типы, объявления функций) \\
\texttt{src/os\_api\_win.c}  & реализация для Windows (Win32 API) \\
\texttt{src/os\_api\_posix.c}& реализация для Linux (POSIX API) \\
\texttt{src/app/parent.c}     & родительский процесс \\
\texttt{src/app/child.c}      & дочерний процесс \\
\texttt{CMakeLists.txt}      & выбор реализации, сборка обоих исполняемых файлов \\
\end{tabular}
\end{center}

\textbf{Основные типы данных:}
\begin{itemize}
  \item \texttt{os\_pipe} --- два дескриптора (\texttt{read}, \texttt{write}) типа
    \texttt{uintptr\_t}; значение \texttt{OS\_INVALID} означает закрытый/несозданный дескриптор;
  \item \texttt{os\_process} --- дескриптор процесса (\texttt{proc}) и его идентификатор
    (\texttt{pid}, используется для логирования).
\end{itemize}

\textbf{Функции единого API и их семантика:}

\begin{center}
\begin{tabular}{ll}
\texttt{os\_self\_pid()}               & PID текущего процесса \\
\texttt{os\_pipe\_create()}           & создание анонимного канала \\
\texttt{os\_pipe\_close\_read/write()} & закрытие конца канала \\
\texttt{os\_pipe\_read()/write()}      & чтение/запись данных через канал \\
\texttt{os\_process\_spawn()}         & запуск дочернего процесса с перенаправлением каналов \\
\texttt{os\_process\_pid()}           & PID запущенного процесса \\
\texttt{os\_process\_wait()}          & ожидание завершения процесса и получение кода выхода \\
\texttt{os\_process\_close()}         & освобождение ресурсов процесса \\
\end{tabular}
\end{center}

Все функции возвращают \texttt{0} при успехе и \texttt{-1} при ошибке
(исключение --- \texttt{os\_pipe\_read}: \texttt{1} -- прочитано ровно \texttt{size} байт,
\texttt{0} -- \texttt{EOF}).

\textbf{Используемые системные вызовы:}

\begin{center}
\begin{tabular}{p{0.4\textwidth}p{0.5\textwidth}}
\textbf{POSIX (Linux)} & \textbf{Win32 (Windows)} \\
\texttt{pipe, getpid, read, write, close} & \texttt{CreatePipe, GetCurrentProcessId, ReadFile, WriteFile, CloseHandle, SetHandleInformation} \\
\texttt{fork, dup2, execv, waitpid}       & \texttt{CreateProcessW, WaitForSingleObject, GetExitCodeProcess} \\
\end{tabular}
\end{center}

\textbf{Обработка ошибок:} каждая платформозависимая функция при сбое формирует диагностическое
сообщение \texttt{report()}: на POSIX используются \texttt{errno} и \texttt{strerror()}, на Windows
--- \texttt{GetLastError()} и \texttt{FormatMessageA()}. Приложение анализирует коды возврата
функций API: любая ошибка приводит к коду возврата процесса \texttt{1}, при этом ресурсы (каналы и
дескриптор процесса) освобождаются всегда. Для предотвращения гибели родительского процесса от
сигнала \texttt{SIGPIPE} при записи в уже закрытый канал в POSIX-реализации сигнал игнорируется, а
ошибка обрабатывается штатно через \texttt{EPIPE}.